%=====================================================================
\chapter{Ensemble Learning}\label{ch:ensembles}
%=====================================================================
\locator{3}{Part III --- Supervised Learning Algorithms}

\begin{outcomes}
By the end of this chapter you will be able to:
\begin{tight}
  \item \textbf{Explain} why a committee of hypotheses can be more accurate than
        its members, and \textbf{compute} the accuracy of a majority vote.
  \item \textbf{Describe} bagging and random forests, and \textbf{explain} which
        part of the error they reduce.
  \item \textbf{Perform} one round of AdaBoost by hand: weighted error, learner
        weight and example reweighting.
  \item \textbf{Describe} gradient boosting and stacking.
  \item \textbf{Choose} between bagging and boosting for a given problem.
\end{tight}
\end{outcomes}

\begin{prereq}
\chref{trees} (decision trees and their instability), \chref{generalization}
(bias and variance) and \chref{bayes} (the Bayes optimal classifier as a weighted
vote of hypotheses).
\end{prereq}

\begin{keyterms}
ensemble $\bullet$ committee $\bullet$ majority (hard) vote $\bullet$ soft vote
$\bullet$ independent errors $\bullet$ bootstrap sample $\bullet$ bagging
$\bullet$ out-of-bag estimate $\bullet$ random forest $\bullet$ feature importance
$\bullet$ boosting $\bullet$ weak learner $\bullet$ AdaBoost $\bullet$ decision
stump $\bullet$ gradient boosting $\bullet$ stacking
\end{keyterms}

\section{Committees of Hypotheses}

The Bayes optimal classifier of \chref{bayes} lets every hypothesis vote, and no
single hypothesis can beat it. Ensemble methods are the practical version of that
idea: train several models, and combine their predictions. They are, in practice,
the most accurate off-the-shelf learners for tabular data --- the random forests
and gradient-boosted trees of this chapter win most such competitions.

\begin{conceptbox}[Why a Committee Can Beat Its Members]
If each member of a committee is right more often than not, and their mistakes are
\emph{independent}, the majority is right more often than any member: for the
majority to be wrong, most members must be wrong on the same example at once.
Everything in this chapter is a way of producing members that are individually
decent and whose errors are as uncorrelated as possible.
\end{conceptbox}

\begin{worked}[The Arithmetic of a Majority Vote]
Five classifiers, each 70\% accurate, with independent errors. The majority is right
when at least three are right:
\[
P(\ge 3 \text{ right}) = \binom{5}{3}0.7^3\,0.3^2 + \binom{5}{4}0.7^4\,0.3 + \binom{5}{5}0.7^5
= 0.309 + 0.360 + 0.168 = \mathbf{0.837}
\]
Twenty-one such classifiers give 0.974; a hundred and one classifiers that are each
only 60\% accurate give 0.979.

\smallskip
\textbf{The catch is ``independent''.} Five copies of the same classifier make the
same mistakes and vote as one: 70\%. Real ensemble members are correlated, so the
gains are smaller than this --- and the whole craft is in reducing the correlation.
\end{worked}

\dsfig{%
\begin{tikzpicture}
\begin{axis}[width=9.4cm, height=5.0cm, axis lines=left,
  xlabel={number of voters (independent errors)}, ylabel={accuracy of the majority},
  tick label style={font=\tiny}, label style={font=\scriptsize},
  xmin=0, xmax=105, ymin=0.5, ymax=1.02,
  legend style={font=\tiny, draw=gray!40, at={(0.98,0.03)}, anchor=south east}]
\addplot[greenhl!70!black, line width=1.2pt, mark=*, mark size=1.2pt] coordinates
  {(1,0.700) (3,0.784) (5,0.837) (7,0.874) (11,0.922) (15,0.950) (21,0.974) (31,0.990) (51,0.999) (75,1.000) (101,1.000)};
\addlegendentry{each voter 70\% accurate}
\addplot[secondary, line width=1.2pt, mark=*, mark size=1.2pt] coordinates
  {(1,0.600) (3,0.648) (5,0.683) (7,0.710) (11,0.753) (15,0.787) (21,0.826) (31,0.872) (51,0.926) (75,0.960) (101,0.979)};
\addlegendentry{each voter 60\% accurate}
\addplot[goldhl, line width=1.2pt, mark=*, mark size=1.2pt] coordinates
  {(1,0.550) (3,0.575) (5,0.593) (7,0.608) (11,0.633) (15,0.654) (21,0.679) (31,0.713) (51,0.764) (75,0.808) (101,0.844)};
\addlegendentry{each voter 55\% accurate}
\end{axis}
\end{tikzpicture}}%
{How majority voting amplifies weak but independent voters. Even 55\%-accurate members
reach 84\% with a hundred voters --- if, and only if, their mistakes are
independent.}{majority}

Voting comes in two forms. A \term{hard vote} counts class predictions. A
\term{soft vote} averages predicted probabilities and picks the largest, which lets
a confident member outweigh several unsure ones and is usually a little better.

\section{Bagging}

\begin{definitionbox}[Bagging (Bootstrap Aggregating)]
From a training set of $n$ examples, draw $B$ \term{bootstrap samples}, each of $n$
examples drawn \emph{with replacement}. Train one model on each sample. Predict by
majority vote (classification) or by averaging (regression).
\end{definitionbox}

Each bootstrap sample leaves some examples out and repeats others: the chance that a
given example is never drawn in $n$ draws is $(1 - 1/n)^n \approx e^{-1} = 0.368$, so
each sample contains about 63.2\% of the distinct examples. The models therefore
differ, and their errors are partly decorrelated.

Bagging helps exactly when the base learner is \emph{unstable}: when small changes
in the training data produce very different models. That describes a fully grown
decision tree (\chref{trees}) precisely --- high variance, low bias --- and bagging
averages the variance away while keeping the low bias. It does little for a stable
learner such as linear regression or $k$-NN with large $k$.

\begin{conceptbox}[A Free Validation Set: the Out-of-Bag Estimate]
Each example is left out of about 37\% of the bootstrap samples. Predict it using
only the models that did not see it, and the resulting \term{out-of-bag} error is an
honest estimate of test error --- computed during training, with no data held back.
\end{conceptbox}

\section{Random Forests}

Bagged trees are still correlated: if one feature is very strong, every tree puts it
at the root and the trees look alike. A \term{random forest} (Breiman, 2001)
decorrelates them further.

\begin{definitionbox}[Random Forest]
Bag fully grown decision trees, with one change: at each split, choose the best test
from a random subset of the features only --- typically $\sqrt{d}$ of the $d$
features for classification. Different trees are forced to use different features,
so their errors are less correlated and the average is better. A forest also
reports each feature's \term{importance}: the total reduction in impurity from splits
on that feature, averaged over the trees.
\end{definitionbox}

Random forests are the safest default for tabular data: few hyperparameters (the
number of trees, which only needs to be large enough, and the feature-subset size),
no scaling, robust to irrelevant features, and hard to overfit by adding trees.

\section{Boosting}

Bagging trains its members independently and in parallel. Boosting trains them in
sequence, each one concentrating on the examples its predecessors got wrong. Its
members are \term{weak learners} --- barely better than chance, typically
\term{decision stumps} (one-split trees) or very shallow trees --- and it is bias,
not variance, that boosting mainly reduces.

\dsfig{%
\begin{tikzpicture}[font=\scriptsize,
  data/.style={rounded corners=2pt, draw=gray!55, fill=gray!8, minimum width=1.5cm,
               minimum height=0.6cm, font=\tiny},
  mdl/.style={rounded corners=3pt, draw=#1, fill=#1!12, text=#1!55!black,
              minimum width=1.2cm, minimum height=0.6cm, font=\tiny\bfseries},
  a/.style={-{Stealth[length=1.8mm]}, gray!60, line width=0.8pt}]
\node[font=\small\bfseries, text=primary, anchor=west] at (-0.6,2.6) {Bagging: in parallel};
\node[data] (d) at (0,1.2) {training data};
\foreach \i/\y in {1/2.0,2/1.2,3/0.4}{
  \node[data] (b\i) at (2.5,\y) {bootstrap \i};
  \node[mdl=secondary] (m\i) at (4.6,\y) {tree \i};
  \draw[a] (d) -- (b\i); \draw[a] (b\i) -- (m\i);}
\node[rounded corners=3pt, draw=greenhl!70!black, fill=greenhl!14, font=\tiny\bfseries,
      minimum width=1.3cm, minimum height=0.8cm, align=center] (v) at (6.6,1.2) {vote /\\average};
\foreach \i in {1,2,3}{\draw[a] (m\i) -- (v);}
\node[font=\tiny, text=gray!55!black, align=center] at (3.5,-0.35) {independent members; reduces variance};
% boosting
\node[font=\small\bfseries, text=accent, anchor=west] at (8.1,2.6) {Boosting: in sequence};
\node[data] (e1) at (8.7,1.2) {weights 1};
\node[mdl=accent] (s1) at (8.7,0.2) {stump 1};
\node[data] (e2) at (10.7,1.2) {weights 2};
\node[mdl=accent] (s2) at (10.7,0.2) {stump 2};
\node[data] (e3) at (12.7,1.2) {weights 3};
\node[mdl=accent] (s3) at (12.7,0.2) {stump 3};
\draw[a] (e1) -- (s1); \draw[a] (e2) -- (s2); \draw[a] (e3) -- (s3);
\draw[a, accent] (s1.east) -- (e2.south west);
\draw[a, accent] (s2.east) -- (e3.south west);
\node[rounded corners=3pt, draw=greenhl!70!black, fill=greenhl!14, font=\tiny\bfseries,
      minimum width=1.3cm, minimum height=0.8cm, align=center] (wv) at (14.7,0.7) {weighted\\vote};
\draw[a] (s3) -- (wv);
\node[font=\tiny, text=gray!55!black, align=center] at (11.6,-0.55) {each member's errors are up-weighted for the next; reduces bias};
\end{tikzpicture}}%
{Bagging trains independent models on resampled data and averages them. Boosting trains
models one after another, re-weighting the data towards the examples still being
misclassified, and combines them with a weighted vote.}{bag-boost}

\begin{definitionbox}[AdaBoost]
Give every training example the weight $1/n$. For rounds $t = 1, \dots, T$:
\begin{tightnum}
  \item Train a weak learner $h_t$ on the weighted data.
  \item Compute its weighted error $\varepsilon_t$ (the total weight of the examples it
        misclassifies) and its say in the final vote,
        $\alpha_t = \tfrac{1}{2}\ln\dfrac{1 - \varepsilon_t}{\varepsilon_t}$.
  \item Multiply the weight of each misclassified example by $e^{\alpha_t}$ and of each
        correct one by $e^{-\alpha_t}$; renormalise to sum to 1.
\end{tightnum}
Predict with the weighted vote $\operatorname{sign}\big(\sum_t \alpha_t h_t(\mathbf{x})\big)$.
\end{definitionbox}

\begin{worked}[One Round of AdaBoost]
Ten examples, each of weight 0.1. The first stump misclassifies three of them.

\smallskip
\textbf{Weighted error:} $\varepsilon_1 = 3 \times 0.1 = 0.3$.

\textbf{Say of the stump:} $\alpha_1 = \tfrac{1}{2}\ln\dfrac{0.7}{0.3} = \tfrac{1}{2}\ln 2.333 = \mathbf{0.424}$.

\textbf{Reweight:} $e^{0.424} = 1.528$ and $e^{-0.424} = 0.655$.
\begin{tight}
  \item three mistakes: $0.1 \times 1.528 = 0.1528$ each (total 0.458);
  \item seven correct: $0.1 \times 0.655 = 0.0655$ each (total 0.458).
\end{tight}
\textbf{Normalise} by the total 0.917: mistakes $\to$ \textbf{0.167} each, correct
$\to$ \textbf{0.071} each.

\smallskip
\textbf{After the update the three mistakes carry exactly half the total weight}
--- always true in AdaBoost. The next stump, trained on these weights, cannot score
better than chance by repeating the first stump's errors, so it is forced to attend
to them. A stump with $\varepsilon = 0.5$ would get $\alpha = 0$: no say at all.
\end{worked}

\textbf{Gradient boosting} generalises the idea to any differentiable loss. Start
with a constant prediction; at each round, fit a small tree to the current
\emph{residuals} (more generally, to the negative gradient of the loss), and add a
fraction $\eta$ --- the learning rate, typically 0.05--0.1 --- of its predictions to
the model. Each tree corrects what the ensemble so far gets wrong. Libraries such as
XGBoost and LightGBM are highly engineered versions, and are the first thing most
practitioners try on a tabular problem.

\dsfig{%
\begin{tikzpicture}
\begin{axis}[width=9.8cm, height=5.2cm, axis lines=left, xmode=log, log basis x=10,
  xlabel={number of trees (log scale)}, ylabel={test error},
  tick label style={font=\tiny}, label style={font=\scriptsize},
  xmin=0.9, xmax=450, ymin=0.05, ymax=0.2,
  xtick={1,5,10,25,50,100,200,400}, xticklabels={1,5,10,25,50,100,200,400},
  yticklabel style={/pgf/number format/fixed}, scaled y ticks=false,
  legend style={font=\tiny, draw=gray!40}, legend pos=outer north east]
\addplot[gray!60, dashed, line width=1pt] coordinates {(0.9,0.129)(450,0.129)};
\addlegendentry{single full tree}
\addplot[secondary, line width=1.2pt, mark=*, mark size=1.2pt] coordinates
  {(1,0.119) (5,0.093) (10,0.083) (25,0.078) (50,0.073) (100,0.073) (200,0.073) (400,0.071)};
\addlegendentry{bagging}
\addplot[greenhl!70!black, line width=1.2pt, mark=square*, mark size=1.2pt] coordinates
  {(1,0.194) (5,0.102) (10,0.077) (25,0.069) (50,0.068) (100,0.071) (200,0.073) (400,0.071)};
\addlegendentry{random forest}
\addplot[accent, line width=1.2pt, mark=triangle*, mark size=1.5pt] coordinates
  {(1,0.157) (5,0.128) (10,0.117) (25,0.097) (50,0.094) (100,0.091) (200,0.089) (400,0.079)};
\addlegendentry{AdaBoost (depth-2 trees)}
\addplot[purplehl, line width=1.2pt, mark=diamond*, mark size=1.5pt] coordinates
  {(1,0.135) (5,0.128) (10,0.114) (25,0.099) (50,0.083) (100,0.081) (200,0.078) (400,0.071)};
\addlegendentry{gradient boosting}
\end{axis}
\end{tikzpicture}}%
{Test error against ensemble size on a 3{,}000-example problem with 5\% label noise.
Every ensemble roughly halves the error of a single tree. Averaging methods level off
quickly; boosting keeps improving for longer.}{ensemble-curves}

\section{Stacking}

\textbf{Stacking} combines \emph{different kinds} of model. Train several base models
--- say logistic regression, an SVM (\chref{svm}) and a random forest --- then train a second-level
model (often logistic regression) whose inputs are the base models' predictions and
whose target is the true label. To avoid leakage, the second-level training data must
be the base models' predictions on examples they did not train on, produced by
cross-validation. Stacking learns how much to trust each base model, and when.

\rowcolors{2}{white}{lightbg}
\begin{longtable}{@{}L{2.8cm}L{3.8cm}L{3.8cm}L{4.4cm}@{}}
\toprule
\rowcolor{primary!15}
\textbf{Method} & \textbf{Members} & \textbf{Mainly reduces} & \textbf{Watch out for} \\
\midrule
\endhead
Bagging & Deep trees on bootstrap samples, in parallel & Variance & Little gain
with stable learners \\
Random forest & Bagged trees with random feature subsets & Variance & Importances
biased towards many-valued features \\
AdaBoost & Stumps or shallow trees, in sequence & Bias (and variance) & Label noise:
mislabelled points get ever larger weights \\
Gradient boosting & Shallow trees fitted to residuals & Bias (and variance) &
Overfits with too many trees at a high learning rate; tune both \\
Stacking & Different model types plus a combiner & Whatever the base models miss &
Leakage if the combiner is trained on in-sample predictions \\
\bottomrule
\end{longtable}

\begin{pitfall}
\begin{tight}
  \item \textbf{Ensembling copies of the same model.} Correlated members vote as
        one; diversity is the whole point.
  \item \textbf{Boosting noisy labels without limit.} AdaBoost's weights chase
        mislabelled points; use fewer rounds, a lower learning rate, or bagging.
  \item \textbf{Reading impurity importances as causal.} They say what the forest
        used, not what drives the outcome, and they favour many-valued features.
  \item \textbf{Training a stacking combiner on in-sample predictions.} Use
        cross-validated predictions.
  \item \textbf{Forgetting the cost.} Hundreds of trees are slower to run and harder
        to explain than one; check that the gain is worth it.
\end{tight}
\end{pitfall}

%---------------------------------------------------------------------
\begin{fourlenses}{Ensembles}
\lensLA{A random forest over engagement features is a strong at-risk model, and
its feature importances give advisers a ranked list of warning signs to discuss.
Present them as ``what the model relies on'', never as ``what causes failure''.}

\lensPM{Gradient-boosted trees are the usual winner for predicting project overrun
or ticket resolution time from mixed tabular data. With a few hundred projects,
keep the trees shallow and the learning rate low, and check with cross-validation
that the ensemble beats a single pruned tree by enough to justify it.}

\lensIOT{Isolation forests --- ensembles of random trees in which anomalies are
isolated in fewer splits --- are a standard unsupervised detector for sensor data.
Forests also parallelise naturally across the cores of an edge gateway.}

\lensSEC{Malware and intrusion classifiers in production are very often
gradient-boosted trees. Ensembles are harder to evade than a single model, because
an attacker must fool many decision paths at once --- harder, not impossible.}
\end{fourlenses}

%---------------------------------------------------------------------
\begin{lab}[One Tree Against Five Ensembles]
\begin{lstlisting}[language=Python]
import numpy as np
from sklearn.datasets import make_classification
from sklearn.model_selection import cross_val_score
from sklearn.tree import DecisionTreeClassifier
from sklearn.linear_model import LogisticRegression
from sklearn.ensemble import (BaggingClassifier, RandomForestClassifier,
                              AdaBoostClassifier, GradientBoostingClassifier,
                              VotingClassifier)

# --- 1. one round of AdaBoost, by hand -----------------------------------
w = np.full(10, 0.1)                            # ten examples, equal weight
wrong = np.array([1, 1, 1, 0, 0, 0, 0, 0, 0, 0], dtype=bool)
eps = w[wrong].sum()                            # weighted error = 0.3
alpha = 0.5 * np.log((1 - eps) / eps)           # say of this learner
w = np.where(wrong, w * np.exp(alpha), w * np.exp(-alpha))
w = w / w.sum()
print(f"alpha = {alpha:.3f}; new weights: wrong {w[0]:.4f}, "
      f"right {w[-1]:.4f}; weight on the mistakes = {w[wrong].sum():.2f}")

# --- 2. one tree against five ensembles ----------------------------------
X, y = make_classification(n_samples=3000, n_features=20, n_informative=8,
                           n_redundant=4, flip_y=0.05, random_state=7)
tree = DecisionTreeClassifier
models = {
    "single tree": tree(random_state=0),
    "bagging (100 trees)": BaggingClassifier(tree(), n_estimators=100,
                                             random_state=0),
    "random forest (100)": RandomForestClassifier(n_estimators=100,
                                                  random_state=0),
    "AdaBoost (depth 2)": AdaBoostClassifier(tree(max_depth=2),
                                             n_estimators=200,
                                             random_state=0),
    "gradient boosting": GradientBoostingClassifier(n_estimators=200,
                                                    random_state=0),
    "vote: LR + tree + RF": VotingClassifier([
        ("lr", LogisticRegression(max_iter=1000)),
        ("dt", tree(max_depth=6, random_state=0)),
        ("rf", RandomForestClassifier(n_estimators=100, random_state=0))],
        voting="soft"),
}
for name, m in models.items():
    s = cross_val_score(m, X, y, cv=5)
    print(f"{name:24s} accuracy {s.mean():.3f} +/- {s.std():.3f}")

# --- 3. a free validation estimate, and what the forest looked at --------
rf = RandomForestClassifier(n_estimators=300, oob_score=True,
                            random_state=0).fit(X, y)
print("random forest out-of-bag accuracy:", round(rf.oob_score_, 3))
top = np.argsort(-rf.feature_importances_)[:5]
print("most used features:", top.tolist(),
      rf.feature_importances_[top].round(3).tolist())
\end{lstlisting}
Every ensemble beats the single tree by four to six points. Now set
\texttt{flip\_y=0.3} (30\% label noise) and compare bagging with AdaBoost again:
which suffers more, and why?
\end{lab}

%---------------------------------------------------------------------
\begin{checkpoint}
\begin{tightnum}
  \item Three independent classifiers are each 80\% accurate. What is the accuracy
        of their majority vote?
  \item A stump has weighted error 0.2. Compute its $\alpha$ and the factors by
        which misclassified and correct examples are re-weighted.
  \item Why does bagging help a decision tree much more than it helps logistic
        regression?
\end{tightnum}
\end{checkpoint}

%---------------------------------------------------------------------
\begin{chaptersummary}
\begin{tight}
  \item A majority of individually decent, independent voters beats every voter;
        the work of ensemble learning is making members diverse.
  \item Bagging trains unstable learners on bootstrap samples and averages them,
        reducing variance; out-of-bag examples give a free error estimate.
  \item Random forests add random feature subsets at each split to decorrelate the
        trees; they are the safest default for tabular data.
  \item Boosting trains weak learners in sequence on re-weighted data; AdaBoost
        gives each learner a say $\alpha = \tfrac{1}{2}\ln\frac{1-\varepsilon}{\varepsilon}$
        and moves half the weight onto its mistakes. It reduces bias but is
        sensitive to label noise.
  \item Gradient boosting fits each new tree to the current residuals; stacking
        learns how to combine different model types.
\end{tight}
\end{chaptersummary}

\begin{reviewq}
  \item \textbf{(MCQ)} A bootstrap sample of $n$ examples contains about what
        fraction of the distinct examples? (a)~50\% (b)~63\% (c)~90\% (d)~100\%.
  \item \textbf{(MCQ)} Random forests differ from bagged trees by: (a)~using stumps
        (b)~re-weighting examples (c)~considering a random subset of features at
        each split (d)~pruning every tree.
  \item \textbf{(MCQ)} In AdaBoost, a weak learner with weighted error 0.5 receives
        a weight $\alpha$ of: (a)~1 (b)~0.5 (c)~0 (d)~$-1$.
  \item \textbf{(Short)} Contrast bagging and boosting in terms of how members are
        trained and which error component each reduces.
  \item \textbf{(Short)} Explain the out-of-bag estimate and why it is honest.
  \item \textbf{(Short)} Why is AdaBoost sensitive to mislabelled training examples?
  \item \textbf{(Applied)} Compute the accuracy of a majority vote of seven
        independent classifiers that are each 65\% accurate.
  \item \textbf{(Applied)} Continue the AdaBoost worked example: the second stump
        misclassifies two of the seven previously correct examples and none of the
        three previous mistakes. Compute $\varepsilon_2$, $\alpha_2$ and the new
        weights.
\end{reviewq}
